

Este trabajo se centra en el estudio y comparación de dos metodologías de valoración de derivados financieros: el modelo de Volatilidad Local Estocástica, en particular el modelo Heston--SLV, y el enfoque de Entropía Relativa Mínima, implementado mediante Weighted Monte Carlo. El objetivo principal es presentar las implementaciones realizadas de ambas metodologías con el fin de obtener una comprensión más amplia de su funcionamiento, sus diferencias y su capacidad para reproducir precios de mercado. \medskip

Tras introducir el marco teórico de valoración neutral al riesgo, el modelo de Black--Scholes y sus principales limitaciones, se presentan ambas metodologías y se desarrolla su implementación numérica. A partir de una superficie de volatilidades implícitas se construyen precios de referencia, se calibra el modelo Heston--SLV y se obtiene una medida posterior mediante MRE/WMC. Finalmente, se comparan los precios obtenidos para opciones europeas vainilla y opciones de tipo \textit{forward-start}.


\vspace{1cm}

\noindent\textbf{\Large Abstract} 

\vspace{1.5cm}

This work focuses on the study and comparison of two methodologies for the valuation of financial derivatives: the Stochastic Local Volatility model, in particular the Heston--SLV model, and the Minimum Relative Entropy approach, implemented through Weighted Monte Carlo. The main objective is to present the implementations carried out for both methodologies in order to obtain a broader understanding of their behaviour, their differences, and their ability to reproduce market prices. \medskip

After introducing the theoretical framework of risk-neutral valuation, the Black--Scholes model and its main limitations, both methodologies are presented and their numerical implementation is developed. Starting from a surface of implied volatilities, reference prices are constructed, the Heston--SLV model is calibrated, and a posterior measure is obtained through MRE/WMC. Finally, the prices obtained for European vanilla options and forward-start options are compared.



